\chapter{Molecular structure and strand direction}
\label{ch:structure}

\section*{A four-letter string is not yet a molecule}
The previous chapter assigned tokens to locations. A molecular encoding needs
more discipline: which physical feature represents a token, and which features
did the representation leave out? We begin with the smallest distinction that
can invalidate an otherwise plausible design. The string \texttt{AACG} and the
string \texttt{TTGC} have complementary letters in corresponding positions,
but writing both from $5'$ to $3'$ does not display an ordinary fully paired
antiparallel duplex.

This chapter builds the vocabulary needed to explain why. You will distinguish
a base from a nucleotide, a covalent backbone from pairing interactions, and
an aligned partner from the same partner rewritten in a standard direction.
The executable example checks these representations. It does not predict
whether strands bind in a particular solution.

\section{The parts and the connections}
\index{nucleotide}\index{strand polarity}
The canonical DNA bases are adenine (A), cytosine (C), guanine (G) and thymine
(T). A base attached to a sugar is a nucleoside; adding phosphate gives a
nucleotide. DNA uses deoxyribose, whereas RNA uses ribose. Along a DNA strand,
phosphodiester linkages join the sugar-phosphate backbone through the $3'$ and
$5'$ positions \citep{cooper2000composition}. The primes refer to positions
on the sugar, not the first and last letters of a displayed string.

These terms describe different levels of an object. A label A in a sequence
usually records base identity while suppressing the associated sugar and
backbone. A drawing can use that shorthand, but its legend should say so.
Otherwise a reader may infer that isolated bases are covalently chained
directly to one another, or that the written characters specify every chemical
feature of the sample.

We will use a precise drawing contract. A solid backbone stroke stands for
covalent connectivity along a strand. A pairing mark between strands stands
for an interaction between bases, not a phosphodiester bond. End labels specify
polarity. The drawings do not specify terminal phosphorylation, protecting
groups or other chemical modifications. A $5'$ label is therefore not a
complete specification for an oligonucleotide order.

\section{Pairing does not erase direction}
In the canonical Watson--Crick pairing model, A pairs with T and C with G.
The two backbones in a conventional duplex are \gls{antiparallel}: tracing one from
$5'$ to $3'$ runs opposite to tracing the other. The familiar double helix has
backbones outside and bases inside \citep{alberts2002dna}. Our flattened view
preserves pairing and polarity but does not depict three-dimensional geometry.
It should not be used to read off angles, distances or helical handedness.

Figure~\ref{fig:polarity} uses the same four positions throughout. Its top row
reads $5'$-AACG-$3'$ from left to right. The directly aligned lower row reads
$3'$-TTGC-$5'$. To report that lower strand in the usual $5'$-to-$3'$ convention,
start at its right-hand end and read leftward. The result is $5'$-CGTT-$3'$.
Nothing about the molecule has changed; we changed the direction in which
we wrote its sequence.

\begin{figure}[p]
\centering
\includegraphics[width=\textwidth]{../build/figures/dna-g06.pdf}
\caption{DNA-G06. An original schematic of canonical pairing and sequence
orientation, using the executable AACG example. Pairing-contact multiplicities
follow the conventional A--T and G--C scheme \citep{alberts2002pairs}.
The backbone and bases are coarse symbols, not atom-resolved structures.
The lower panel changes written orientation, not chemical identity.}
\label{fig:polarity}
\end{figure}

\section{Two transformations with different meanings}
Let $s=s_1\cdots s_n$ be a nonempty canonical sequence written from $5'$ to
$3'$. Define the symbol involution $c$ by
\[
 c(A)=T,\quad c(T)=A,\quad c(C)=G,\quad c(G)=C.
\]
Extending $c$ position by position gives $c(s)=c(s_1)\cdots c(s_n)$.
Let $R(s)=s_n\cdots s_1$. The \gls{reverse-complement} is\index{reverse complement}
\[
 \operatorname{RC}(s)=R(c(s))=c(s_n)\cdots c(s_1).
\]
These are mathematical definitions. Their interpretation as the two written
views in the figure depends on the canonical, fully paired, antiparallel
model just specified. The functions alone do not establish a physical duplex.

\begin{proposition}[Orientation round trip]
On canonical sequences, $c$ and $R$ commute and are both involutions.
Consequently $\operatorname{RC}(\operatorname{RC}(s))=s$.
\end{proposition}
\begin{proof}
At position $i$, both $R(c(s))$ and $c(R(s))$ contain $c(s_{n+1-i})$.
Thus the transformations commute. Applying $c$ twice restores each base;
applying $R$ twice restores each position. Hence
$(Rc)(Rc)=RRcc$ is the identity.
\end{proof}

An involution test is necessary but not sufficient to identify the intended
transformation. Reversal alone, basewise complement alone, and even the identity
function each pass a round-trip test. An asymmetric example provides the
missing discrimination. For AACG, these outputs are respectively GCAA, TTGC
and AACG, whereas the required reverse complement is CGTT. A test suite should
check both the algebraic property and at least one independently specified
orientation example.

\section{A representation that exposes its assumptions}
Our reference distinguishes a stored \texttt{Strand}, always written $5'$ to
$3'$, from a \texttt{SequenceView}, whose direction is part of its data.
A \texttt{Duplex} stores both member strands in standard orientation and
validates that the second is the reverse complement of the first. Displaying
the second strand beneath the first reverses its written order and labels it
$3'$ to $5'$.

\begin{minipage}{\textwidth}
\begin{Verbatim}[fontsize=\small]
from minidna.duplex import Duplex
duplex = Duplex.from_sequence("AACG")
for row in duplex.rows():
    print(row)
print(duplex.bottom)
\end{Verbatim}
\end{minipage}

The records printed by this example are summarized in
Table~\ref{tab:orientation}. Both rows of the paired view contain four
nucleotides. The ideal duplex contains four base pairs, not eight base pairs;
counting nucleotide residues across both strands would give eight.

\begin{table}[htbp]
\centering
% Generated by scripts/chapter02_artifacts.py; do not edit.
\begin{tabular}{lll}
\toprule View & Direction & Sequence \\
\midrule
Top & 5' to 3' & \texttt{AACG} \\
Aligned partner & 3' to 5' & \texttt{TTGC} \\
Partner, standard & 5' to 3' & \texttt{CGTT} \\
\bottomrule\end{tabular}

\caption{Generated views of the same ideal duplex. Direction belongs to the
representation and must travel with the letters.}
\label{tab:orientation}
\end{table}

This API excludes empty strands, ambiguous symbols, mismatches and overhangs.
It normalizes case and removes whitespace, including Unicode whitespace.
It does not parse decorated strings such as \texttt{5'-AACG-3'} as sequence
input; those end labels belong to the view. A rejected input is not a claim
that the corresponding chemistry is impossible. It says the input falls
outside this deliberately small model.

The earlier \texttt{Strand.complement()} method is retained for compatibility
as a symbol substitution reinterpreted as a new standard-orientation sequence.
It must not be used as a duplex constructor. The new
\texttt{aligned\_complement()} method returns a direction-labeled view; the
existing \texttt{reverse\_complement()} returns the partner in standard
orientation. Documenting this distinction resolves the previous teaching
ambiguity without silently changing existing callers' results.

The tests enumerate all 1364 canonical sequences of lengths one through five.
They check every displayed pairing column and both directions, then verify
the round trip. A separate test rejects the incorrect pair of standard-form
strands AACG and TTGC. Passing these tests certifies representation invariants,
not affinity, reaction yield or biochemical selectivity.

\section{Direction becomes a process constraint}
For ordinary template-directed DNA primer extension, polymerase adds an
incoming nucleotide at the primer's $3'$ hydroxyl end. New DNA grows in the
$5'$-to-$3'$ direction while the paired template is traversed oppositely
\citep{cooper2000replication}. These labels describe different strands and
must never be combined into a statement that the new strand grows $3'$ to $5'$.

Consider the purely symbolic alignment of a template written
$3'$-TTGC-$5'$ and a primer written $5'$-AA-$3'$ beneath its first two
positions. Continuing the canonical sequence would next add C and then G at
the primer's $3'$ end, yielding AACG. Rewriting the template as CGTT in
standard orientation does not change that answer; it changes how one locates
the primer relative to the displayed template. This exercise specifies
orientation only. A two-base primer here is a pedagogical token, not a
laboratory design or a claim that extension would work under real conditions.

The same care is needed when a drawing is rotated. A left-pointing synthesis
arrow can still represent $5'$-to-$3'$ growth if the strand's $5'$ end lies to
the right. Page coordinates and chemical polarity are different coordinate
systems. Inspect the end labels before deciding that a leftward arrow is wrong.

\section{What sequence logic leaves unresolved}
A completely complementary sequence pair is an input to a physical question,
not its answer. To estimate binding under specified conditions, a later model
would need more than these letters and an equality predicate. We would have
to declare which physical states are represented, the relevant environment,
the measured quantity and how predictions are checked against observations.
The present API intentionally has no method called ``binding probability.''

Likewise, a mismatch is a difference from the canonical paired model, not an
automatic impossibility declaration. A sequence validator should not turn a
restricted alphabet into an assertion that modified bases do not exist. These
boundaries matter because a program can be perfectly correct for its own
domain while being an unsuitable model for the experimental question.

\section{Exercises and solution notes}
\begin{enumerate}
\item Write the aligned partner of $5'$-AGTC-$3'$ with both end labels. Then
write that partner in standard orientation.
\item Why can testing only ATAT conceal a missing reversal in some operations?
Compare identity and reverse complement on this input.
\item Give three wrong implementations that pass the reverse-complement
round-trip test. Which additional test rejects them for AACG?
\item An ideal duplex has two fully paired strands of length 12. How many
base pairs and nucleotide residues does it contain?
\item In the symbolic primer example, identify the growing end, the next base
and the direction in which the template is traversed.
\item State two experimental conclusions that cannot be obtained by constructing
our ideal duplex from the sequence AACG.
\end{enumerate}

\paragraph{Solution notes.}
(1) The aligned row is $3'$-TCAG-$5'$; the standard-form partner is
$5'$-GACT-$3'$. (2) ATAT equals its reverse complement, so identity can masquerade
as RC on this example. (3) Identity, reversal and basewise complement are all
involutions; a direct expected-value test requiring CGTT rejects all three.
(4) Twelve base pairs and 24 nucleotide residues. (5) The primer's $3'$ end
grows; C is next; the template is traversed $3'$ to $5'$. (6) For example,
neither binding probability nor successful polymerase extension is established.

\section*{The boundary we can now maintain}
We can represent canonical DNA sequences without losing their direction and
distinguish an exact string identity from a physical prediction. The next
chapter develops molecular recognition and its energetic limits, where the
gap between compatible letters and a successful interaction becomes central.
